% ParserLab — Presentación de la implementación
% Compilar: pdflatex presentacion.tex (necesita beamer + metropolis)
% Si no tienes metropolis: cambia \usetheme{metropolis} por \usetheme{Madrid}
%                          y comenta los \metroset si los hay.

\documentclass[aspectratio=169,10pt]{beamer}

\usepackage[utf8]{inputenc}
\usepackage[T1]{fontenc}
\usepackage[spanish]{babel}
\usepackage{xcolor}
\usepackage{tikz}
\usepackage{booktabs}
\usepackage{listings}

\usetheme{metropolis}
\metroset{block=fill,progressbar=frametitle}

% Paleta consistente con la app
\definecolor{plPrimary}{HTML}{2563EB}
\definecolor{plSuccess}{HTML}{10B981}
\definecolor{plWarning}{HTML}{F59E0B}
\definecolor{plDanger}{HTML}{EF4444}
\definecolor{plInk}{HTML}{0F172A}
\setbeamercolor{progress bar}{fg=plPrimary}
\setbeamercolor{frametitle}{bg=plPrimary,fg=white}
\setbeamercolor{alerted text}{fg=plDanger}

\lstdefinestyle{gram}{
  basicstyle=\ttfamily\footnotesize,
  keywordstyle=\color{plPrimary},
  showstringspaces=false,
  frame=single,
  rulecolor=\color{gray!30},
  framesep=4pt,
  backgroundcolor=\color{gray!5},
}

\title{\textbf{ParserLab}}
\subtitle{Analizadores sintácticos interactivos: \\
descenso recursivo, LL(1), LR(0), SLR(1), LALR(1), LR(1)}
\author{Compiladores — Examen 1 (Puntos extra)}
\institute{UTEC · 2026-1}
\date{18 de mayo de 2026}

\begin{document}

% =============== Frame 1 — Portada ===============
\maketitle

% =============== Frame 2 — Objetivo ===============
\begin{frame}{Objetivo y motivación}
  \begin{block}{¿Qué es ParserLab?}
    Aplicación web pedagógica que ejecuta y \alert{visualiza} los
    seis algoritmos clásicos de análisis sintáctico sobre la misma
    gramática y entrada.
  \end{block}

  \begin{itemize}
    \item Una sola UI: el usuario \textbf{compara} cómo cada parser
      reacciona ante la misma gramática.
    \item Cada análisis muestra \textbf{traza paso a paso}, \textbf{construcción
      de tablas} y \textbf{validación de cadenas}.
    \item Integra \textbf{IA heurística} para explicar errores y sugerir
      transformaciones LL(1).
    \item Pensada para enseñar la \textbf{jerarquía progresiva} de poder
      expresivo (LR(0) $\subset$ SLR(1) $\subset$ LALR(1) $\subset$ LR(1)).
  \end{itemize}
\end{frame}

% =============== Frame 3 — Jerarquía de parsers ===============
\begin{frame}{Jerarquía implementada}
  \begin{columns}[T,onlytextwidth]
    \column{0.48\textwidth}
    \begin{block}{\textcolor{plPrimary}{Top-Down}}
      \begin{itemize}
        \item Descenso recursivo predictivo
        \item LL(1) con tabla M[A,a]
      \end{itemize}
      \vspace{6pt}
      Requieren: \emph{sin recursión izquierda directa} y
      \emph{factorización de prefijos}.
    \end{block}

    \column{0.48\textwidth}
    \begin{block}{\textcolor{plWarning}{Bottom-Up}}
      \begin{itemize}
        \item LR(0) — sin lookahead
        \item SLR(1) — reduce restringido a FOLLOW
        \item LALR(1) — fusión de núcleos LR(1)
        \item LR(1) canónico
      \end{itemize}
      Cada uno resuelve más conflictos que el anterior.
    \end{block}
  \end{columns}

  \vspace{8pt}
  \centering
  \footnotesize
  \textbf{6 algoritmos · 1 backend Flask · 1 UI React}
\end{frame}

% =============== Frame 4 — Arquitectura ===============
\begin{frame}{Arquitectura}
  \centering
  \begin{tikzpicture}[
    node distance=10mm,
    box/.style={draw,rounded corners,minimum height=10mm,minimum width=28mm,
                font=\small,align=center,fill=plPrimary!10,draw=plPrimary},
    sub/.style={draw,rounded corners,minimum height=6mm,minimum width=20mm,
                font=\scriptsize,align=center,fill=gray!5,draw=gray!50},
    arr/.style={->,>=stealth,thick,plInk}
  ]
    \node[box] (front) {Frontend\\ React + Bootstrap};
    \node[box,right=20mm of front] (back) {Backend\\ Flask + Python};
    \node[box,right=20mm of back] (viz) {Renderer\\ @viz-js/viz (WASM)};

    \node[sub,below=of front,xshift=-18mm] (ui1) {ParserForm};
    \node[sub,below=of front] (ui2) {ParserResult};
    \node[sub,below=of front,xshift=18mm] (ui3) {Compare};

    \node[sub,below=of back,xshift=-15mm] (b1) {6 runners};
    \node[sub,below=of back,xshift=10mm] (b2) {parser\_insights\\ (IA)};
    \node[sub,below=of b1,xshift=8mm] (b3) {nfa\_lr\\ dot\_export};

    \draw[arr] (front) -- node[above,font=\scriptsize]{POST /api/parse} (back);
    \draw[arr] (back.south west) |- ([yshift=2mm]b1.north);
    \draw[arr] (back.south) -- (b2.north);
    \draw[arr,plPrimary] (front) -- node[above,font=\scriptsize]{dot string} (viz);
  \end{tikzpicture}

  \vspace{4pt}
  \footnotesize
  \textbf{POST /api/parse} con \texttt{\{grammar,input,parser\}} →
  payload único con traza, tablas, autómatas (DOT), parse tree, asistente IA.
\end{frame}

% =============== Frame 5 — Lo que ve el usuario ===============
\begin{frame}{Por cada parser: traza · tablas · veredicto}
  \begin{itemize}
    \item \textbf{Simulación paso a paso} con stepper (Inicio/Anterior/Siguiente/Fin)
      que sincroniza con la tabla resaltando la fila actual.
    \item \textbf{Tablas ACTION y GOTO} con color-coding semántico:
      \textcolor{plSuccess}{shift}, \textcolor{plPrimary}{reduce},
      \textcolor{plWarning}{accept}, \textcolor{plDanger}{error}.
    \item \textbf{FIRST / FOLLOW / tabla M[A,a]} para top-down; lista de
      conflictos detectados.
    \item \textbf{Doble veredicto}: si la \emph{cadena} fue aceptada y si la
      \emph{gramática} pertenece a la clase del parser.
  \end{itemize}

  \begin{exampleblock}{Ejemplo de mensaje doble}
    \textcolor{plDanger}{$\times$ Entrada rechazada} ·
    \textcolor{plWarning}{$\triangle$ Gramática NO es LR(0): 4 conflictos reduce-reduce.}
  \end{exampleblock}
\end{frame}

% =============== Frame 6 — Visualizaciones ===============
\begin{frame}{Visualizaciones avanzadas}
  Construidas con \textbf{Graphviz} (vía \texttt{@viz-js/viz}, WebAssembly):
  layout automático \texttt{dot}, zoom/pan, interactividad inyectada en el SVG.

  \begin{columns}[T]
    \column{0.5\textwidth}
    \begin{block}{Bottom-Up}
      \begin{itemize}
        \item \textbf{AFN de ítems} LR(0) o LR(1): cada ítem es un nodo,
          $\varepsilon$-aristas punteadas son el \emph{closure}.
        \item \textbf{DFA canónico}: resultado de subset-construction.
        \item Click en nodo $\to$ panel con los ítems del estado.
      \end{itemize}
    \end{block}

    \column{0.5\textwidth}
    \begin{block}{Comunes}
      \begin{itemize}
        \item \textbf{Parse tree} / derivación renderizado con Graphviz.
        \item \textbf{Teclado virtual} de símbolos formales
          ($\varepsilon, \to, \cup, \cap, \forall, \exists, \Sigma, \ldots$).
        \item \textbf{Export} a \texttt{.dot}, \texttt{.png} y PDF de tablas.
      \end{itemize}
    \end{block}
  \end{columns}
\end{frame}

% =============== Frame 7 — IA / Asistente ===============
\begin{frame}{Asistente inteligente}
  Módulo \texttt{parser\_insights.py} — heurístico, sin llamadas externas
  (reproducible sin API keys).

  \begin{itemize}
    \item \textbf{Explicación de errores en lenguaje natural}: identifica
      el último paso de error en la traza y describe la causa probable
      (token inesperado, lookahead sin acción, falta de GOTO, $\ldots$).
    \item \textbf{Recomendaciones ante ambigüedad}: lista los conflictos
      shift-reduce / reduce-reduce y sugiere el siguiente nivel
      (LR(0) $\to$ SLR(1) $\to$ LR(1)).
    \item \textbf{Sugerencias hacia LL(1)}: detecta recursión izquierda
      directa y prefijos comunes, propone reescrituras
      (factorización, eliminación de recursión).
    \item \textbf{Linter en tiempo real} en el editor de gramática:
      errores estructurales, no-terminales sin uso, producciones duplicadas.
  \end{itemize}
\end{frame}

% =============== Frame 8 — Funcionalidades innovadoras ===============
\begin{frame}{Funcionalidades de valor agregado}
  \begin{columns}[T,onlytextwidth]
    \column{0.48\textwidth}
    \begin{block}{UX y productividad}
      \begin{itemize}
        \item \textbf{Sidebar} con jerarquía top-down/bottom-up
        \item \textbf{Modo comparación} side-by-side con bar de diferencias
        \item \textbf{Historial} persistente en \texttt{localStorage} (20 entradas, dedup)
        \item \textbf{Catálogo de ejemplos} filtrados por parser, con tag $\triangle$ cuando muestran conflicto
        \item \textbf{Light / dark theme} con tokens CSS
      \end{itemize}
    \end{block}

    \column{0.48\textwidth}
    \begin{block}{Despliegue \& export}
      \begin{itemize}
        \item \textbf{PWA} instalable (manifest, standalone)
        \item \textbf{Docker compose} (un comando levanta API + Web)
        \item Deploy público en Render
        \item Export tablas a \textbf{PDF}, grafos a \textbf{.dot} y \textbf{.png}
      \end{itemize}
    \end{block}
  \end{columns}

  \vspace{4pt}
  \centering\footnotesize
  Diseño system con paleta semántica · tipografía Inter + JetBrains Mono.
\end{frame}

% =============== Frame 9 — Suite de pruebas ===============
\begin{frame}{Suite de pruebas multi-parser}
  \texttt{python run\_all\_tests.py} cubre los 6 parsers contra 13 casos:
  \textbf{78 corridas}, 0 WRONG, 0 ERROR.

  \begin{center}
  \footnotesize
  \begin{tabular}{lrrrrr}
    \toprule
    \textbf{Parser} & \textbf{OK} & \textbf{CONF-WRONG} & \textbf{WRONG} & \textbf{ERROR} \\
    \midrule
    RD       &  9 & 4 & 0 & 0 \\
    LL(1)    &  9 & 4 & 0 & 0 \\
    LR(0)    &  7 & 6 & 0 & 0 \\
    SLR(1)   & 10 & 3 & 0 & 0 \\
    LALR(1)  & 11 & 2 & 0 & 0 \\
    LR(1)    & \textbf{13} & 0 & 0 & 0 \\
    \bottomrule
  \end{tabular}
  \end{center}

  \begin{exampleblock}{Lectura pedagógica}
    La jerarquía LR(0) $\to$ SLR(1) $\to$ LALR(1) $\to$ LR(1) se observa
    directamente en el conteo de OK. CONF-WRONG indica que la gramática
    cae fuera de la clase del parser, no un bug.
  \end{exampleblock}
\end{frame}

% =============== Frame 10 — Cierre / demo ===============
\begin{frame}{Demo \& despliegue}
  \begin{block}{Stack}
    React 19 · Flask · @viz-js/viz · svg-pan-zoom · jsPDF · Bootstrap 5
  \end{block}

  \begin{block}{Levantar localmente}
    \texttt{docker compose up --build}   \quad → \quad
    Web en \texttt{http://localhost:8080}
  \end{block}

  \begin{block}{Uso de IA en el desarrollo}
    Diseño, refactor y depuración asistidos por Claude · IA integrada
    en producto vía \texttt{parser\_insights} (heurística determinista).
  \end{block}

  \vspace{6pt}
  \centering
  \Large\textbf{¡Gracias!}\\
  \small
  Preguntas y demo en vivo.
\end{frame}

\end{document}
